\section{Prototype Evaluation}

\subsection{RAG Variant Comparison}
Prototype evaluation was performed using benchmark feedback cases and several
retrieval configurations:
\begin{itemize}
    \item \texttt{no\_rag};
    \item \texttt{direct\_k8};
    \item \texttt{direct\_k8\_unfiltered};
    \item \texttt{hyde\_k8};
    \item \texttt{hyde\_k8\_unfiltered}.
\end{itemize}

The outputs from \texttt{notebooks/rag\_debug\_demo.ipynb} illustrate how
retrieval affects score calibration and feedback interpretation.

The comparison is designed to answer a practical calibration question: does
retrieval help the system place feedback examples closer to expected quality
bands? The no-RAG variant provides a prompt-only baseline. Direct retrieval
tests whether evidence from the knowledge base improves evaluation when the
retrieval query is based directly on the transcript and the feedback-quality
criteria. HyDE retrieval adds one additional step: the LLM first generates a
short hypothetical guidance passage that resembles the kind of educational
evidence that would be useful for judging the transcript. This generated passage
is then embedded and used to retrieve knowledge-base chunks. In other words,
direct retrieval searches with the evaluation query itself, whereas HyDE
retrieval searches with a generated document-like representation of the
evidence that the system hopes to find.

The comparison also separates filtered and unfiltered retrieval. In the
unfiltered variants, criterion-aware query construction and quality reranking
are disabled, so retrieval remains closer to the raw transcript-based vector
similarity ranking:

\begin{center}
\texttt{query -> retrieve candidates -> sort by relevance -> take final\_k chunks}
\end{center}

The filtered variants add two refinements: the retrieval query is expanded with
the feedback-quality criteria, and post-retrieval quality control is applied
before the final context is selected:

\begin{center}
\texttt{query -> retrieve candidates -> quality reranking -> remove low-value chunks -> take final\_k chunks}
\end{center}

This distinction is important because vector similarity alone can retrieve
chunks that are semantically related but not educationally useful, such as
generic introductions, conclusions, reference sections, or noisy document
fragments. The criterion-aware query pushes retrieval toward guidance about the
six rubric dimensions, while quality reranking gives higher priority to chunks
that mention specific observed behavior, objective wording, changeable
improvement areas, and concrete next steps. The unfiltered variants therefore
answer the question of how far raw retrieval can go, while the filtered variants
test whether task-specific query construction and evidence-selection logic
improve calibration.

The benchmark cases are grouped by expected quality bands such as weak,
mediocre, strong, and edge cases. This grouping is useful because a model may
perform well on obviously strong feedback while still struggling with borderline
or moderately useful feedback. For formative examiner training, those middle and
edge cases are especially important because they resemble realistic feedback
that is partially useful but still improvable.

\begin{table}[H]
    \centering
    \small
    \begin{tabular}{L{0.25\linewidth}L{0.20\linewidth}L{0.15\linewidth}L{0.28\linewidth}}
        \toprule
        Variant & Avg. distance to gold band & Within band & Interpretation \\
        \midrule
        \texttt{hyde\_k8} & 2.19 & 10/16 & Best overall alignment \\
        \texttt{hyde\_k8\_unfiltered} & 2.31 & 9/16 & Very close to filtered HyDE \\
        \texttt{direct\_k8} & 2.56 & 8/16 & Useful, but weaker than HyDE \\
        \texttt{direct\_k8\_unfiltered} & 3.00 & 6/16 & Lower calibration \\
        \texttt{no\_rag} & 3.31 & 7/16 & Prompt-only baseline \\
        \bottomrule
    \end{tabular}
    \caption{Illustrative RAG comparison from the debug notebook. Lower
    distance to the gold band is better.}
    \label{tab:rag-variant-comparison}
\end{table}

\subsection{LLM-as-a-Judge Validation}
In addition to the deterministic distance-to-gold-band comparison, an
LLM-as-a-judge step was used as a secondary calibration check. The judge read
the already generated evaluation outputs from the RAG comparison and compared
them with the benchmark expectations. In the debug notebook, this produced 80
completed judgments, corresponding to 16 benchmark cases evaluated under five
variants.

The judge was not treated as ground truth. Instead, it acted as a structured
second lens beside the numeric band-distance calculation. This is important
because an overall score can be numerically close to the target band while the
reasoning remains educationally weak, or a criterion score can appear plausible
while the explanation does not match the intended benchmark expectation. The
judge therefore assessed four complementary dimensions:
\begin{itemize}
    \item \textbf{Overall band fit}: whether the final feedback-quality score
    places the examiner's feedback in the expected benchmark category, such as
    weak, mediocre, strong, or edge-case feedback. This checks the global
    calibration of the evaluation.
    \item \textbf{Criterion band fit}: whether the six rubric criteria are
    scored consistently with the benchmark expectation. For example, feedback
    that contains no concrete next step should not receive a high score for the
    improvement-plan criterion, even if its general tone is acceptable.
    \item \textbf{Benchmark expectation match}: whether the generated summary,
    reasoning, and suggestions explain the same strengths and weaknesses that
    the benchmark case was designed to test. This checks whether the system
    understood why a feedback example should be considered weak, partial, or
    strong.
    \item \textbf{Directional correctness}: whether the evaluation moves in the
    correct formative direction. In this project, that means avoiding inflated
    scores for vague feedback, avoiding excessive penalties for clearly
    structured feedback, and preserving the intended order between weak,
    mediocre, strong, and ambiguous examples.
\end{itemize}

\begin{table}[H]
    \centering
    \scriptsize
    \setlength{\tabcolsep}{3pt}
    \begin{tabular}{L{0.22\linewidth}L{0.12\linewidth}L{0.13\linewidth}L{0.13\linewidth}L{0.15\linewidth}L{0.14\linewidth}}
        \toprule
        Variant & Mean judge score & Overall band fit & Criterion band fit & Benchmark match & Directional correctness \\
        \midrule
        \texttt{hyde\_k8} & 62.97 & 71.25 & 52.50 & 51.56 & 76.56 \\
        \texttt{hyde\_k8\_unfiltered} & 62.61 & 67.19 & 55.88 & 54.88 & 72.50 \\
        \texttt{direct\_k8} & 55.88 & 67.50 & 46.94 & 42.50 & 66.56 \\
        \texttt{direct\_k8\_unfiltered} & 54.61 & 53.12 & 49.81 & 49.88 & 65.62 \\
        \texttt{no\_rag} & 52.70 & 54.38 & 48.31 & 47.19 & 60.94 \\
        \bottomrule
    \end{tabular}
    \caption{LLM-as-a-judge summary from the debug notebook. Scores are
    reported on a 0--100 scale; higher values indicate stronger agreement with
    the benchmark expectation.}
    \label{tab:llm-judge-comparison}
\end{table}

The judge results support the same overall direction as the deterministic
comparison in Table~\ref{tab:rag-variant-comparison}. The strongest mean judge
score belongs to the filtered HyDE variant, with the unfiltered HyDE variant
following very closely. This indicates that the advantage of HyDE is not only
visible in the numeric distance to the gold band, but also in the judged quality
of the generated evaluation. The filtered HyDE variant achieved the highest
overall band fit and directional correctness, suggesting that it was best at
placing examples in the right broad quality region and preserving the expected
weak/mediocre/strong/edge direction.

The unfiltered HyDE variant was very close and performed slightly better on
criterion band fit and benchmark expectation match. This suggests that raw HyDE
retrieval may sometimes preserve criterion-level details that are useful for
the judge. However, the filtered \texttt{hyde\_k8} variant remained stronger as
an operational default because it gave the best balance between overall
calibration, directional correctness, and robustness against low-value context.

\subsection{Interpretation of Results}
Taken together, the deterministic and judge-based results show a consistent
pattern. Retrieval improves the evaluation when it provides evidence that is
well aligned with the feedback-quality rubric, and HyDE appears to improve this
alignment more reliably than direct retrieval alone. The selected
\texttt{hyde\_k8} configuration achieved the lowest average distance to the
gold band, the highest number of in-band cases, the highest mean judge score,
and the strongest directional correctness. This made it the most convincing
default configuration for the current prototype.

The comparison also shows that the no-RAG baseline is not useless. A carefully
designed rubric prompt already gives the LLM meaningful structure, and this is
why the prompt-only baseline can still produce reasonable evaluations in some
cases. The added value of RAG is therefore not that it creates the rubric from
nothing, but that it improves grounding and calibration by placing the
transcript beside relevant educational evidence.

The difference between filtered and unfiltered retrieval should be interpreted
carefully. Filtering and quality reranking help remove generic or noisy chunks
and prioritize evidence about behavior-specific, objective, and actionable
feedback. At the same time, the close performance of
\texttt{hyde\_k8\_unfiltered} shows that the generated HyDE passage itself
already retrieves useful evidence. This means that both components matter: HyDE
improves the semantic match between the evaluation task and the knowledge base,
while filtering tries to make the final context more reliable.

The most difficult category remains moderately useful feedback. Strong examples
are easier to recognize because they satisfy most criteria, and weak examples
are easier because many criteria are absent. Mediocre feedback often contains a
mixture of useful and missing elements, making calibration more sensitive to
retrieved evidence, scoring anchors, and prompt wording. This observation is
important for future validation because real examiner feedback is likely to fall
frequently into this intermediate range. Future work should therefore repeat
the comparison with a larger benchmark set and with the currently available
KISSKI evaluation model to test whether the same ranking remains stable.
